\section{Borrowing GPUs}
\label{sec:borrowing}

\subsection{The verbs}

\begin{lstlisting}
gpu status        # who holds what, what is free
gpu claim 2       # take 2 (all-or-nothing: 2 or none)
gpu shell         # shell inside your container
gpu shell CMD     # run one command in it, non-interactively
gpu release       # give them back
\end{lstlisting}

\cmd{gpu} on the developer VM is a thin relay: it re-quotes its arguments and
forwards them over SSH to \cmd{camimo}, where the real tool runs. Because it
requests a remote TTY only when it has one, it works from scripts and from
automated callers as well as interactively.

\subsection{What a claim actually creates}
\cmd{gpu claim} allocates cards in pool state and then starts a container
pinned to exactly those cards:

\begin{lstlisting}
docker run -d --name gpu-eugen \
  --gpus "device=1" \
  --label gpu.owner=eugen --label gpu.devices=1 \
  --restart no \
  --user 3003:3000 \
  -e HOME=/home/eugen \
  -v /shared:/shared \
  --shm-size 16g \
  nvidia/cuda:12.6.0-base-ubuntu24.04 sleep infinity
\end{lstlisting}

Pool state lives in \cmd{/var/lib/gpu-pool/state.json} under a \cmd{flock}, and
every command reconciles that file against \cmd{docker ps} before reporting or
allocating. A container that dies therefore returns its cards on its own. If
the state file is unreadable the tool refuses to write rather than falling back
to ``the pool is empty'', which would hand out cards already in use.

\subsection{Why a container at all}

\begin{table}[!ht]
  \centering
  \renewcommand{\arraystretch}{1.25}
  \begin{tabularx}{\linewidth}{@{}p{4.2cm} X@{}}
    \toprule
    \textbf{Property} & \textbf{Consequence} \\
    \midrule
    \cmd{-{}-gpus "device=N"} &
    Isolation is structural, not by convention. Other cards are invisible;
    yours renumbers to 0. A stray \cmd{CUDA\_VISIBLE\_DEVICES} cannot reach a
    colleague's run. \\
    Container is the unit of accounting &
    \cmd{gpu release} is \cmd{docker rm -f}; reconciliation against
    \cmd{docker ps} frees the cards. \\
    No reaper, by design &
    A claim is released by its owner, never by a timer --- so a long training
    run is safe, and a forgotten claim blocks the pool. \\
    All-or-nothing allocation &
    Asking for 4 with 3 free yields nothing, rather than a half-sized
    allocation that silently invalidates the run. \\
    \cmd{-{}-user uid:3000} &
    Files created on \cmd{/shared} are owned by you in the \cmd{team} group,
    not by root. \\
    \bottomrule
  \end{tabularx}
  \caption{What the container arrangement buys.}
  \label{tab:gpu-container-properties}
\end{table}

\subsection{Consequences to plan around}
The image is \cmd{nvidia/cuda:12.6.0-base-ubuntu24.04}. The \emph{base} variant
carries the CUDA runtime and nothing else.

\begin{table}[!ht]
  \centering
  \renewcommand{\arraystretch}{1.25}
  \begin{tabularx}{\linewidth}{@{}p{5.1cm} X@{}}
    \toprule
    \textbf{Constraint} & \textbf{Why} \\
    \midrule
    No \cmd{python3}, \cmd{curl}, \cmd{git}, \cmd{tmux} & \cmd{-base} image variant \\
    No \cmd{sudo}; \cmd{apt install} fails & runs as your uid, not root \\
    Service ports unreachable from outside & default docker bridge, isolated netns \\
    Filesystem is ephemeral & \cmd{gpu release} deletes the container \\
    \bottomrule
  \end{tabularx}
  \caption{Constraints imposed by the claim container.}
  \label{tab:gpu-constraints}
\end{table}

Anything that must run inside the container therefore has to be a
self-contained binary staged in from outside.
